\section{Discussion}
\label{sec:discussion}

\subsection{Connectivity change and comparison with previous work}
\TODO{PLAIN-LANGUAGE GUIDE: begin with the simplest overall result: the modeled ease
of moving between cores changed only slightly overall, but some individual links
changed more than others. Immediately state that the model describes landscape
structure and does not show where cheetahs actually moved.}

% Interpret individual increases and decreases in this same subsection.
\TODO{PLAIN-LANGUAGE GUIDE: describe where modeled connections became easier or
harder and which landscape changes occurred in the same places. Say "occurred in the
same places," not "caused." The model found mixed change rather than a broad decline.
Do not describe connectivity by country because core pairs were not assigned to
individual countries.}

% Compare these findings with previous regional studies here.
\TODO{PLAIN-LANGUAGE GUIDE: compare the general mapping approach with the elephant and
multi-species studies \parencite{naidoo2024kaza,brennan2020multispecies} and with the
fence study \parencite{osipova2018fencing}. Explain that similar mapping tools can be
used for different species, but the numerical resistance values cannot simply be
copied because each species responds differently to roads, vegetation, fences, and
people.}


\subsection{Protected areas and conservation context}
\TODO{PLAIN-LANGUAGE GUIDE: the strongest modeled concentrations often pass through
protected areas, but the wider land connecting those areas is less protected. Explain
that conservation therefore also depends on farms, communal lands, and other land
outside formal parks. This does not prove that protected areas failed. Also note that
parks are often located on land with lower agricultural value, which may partly shape
the pattern.}

\subsection{Priority links and field-survey needs}
\TODO{PLAIN-LANGUAGE GUIDE: links that stayed important under several assumptions are
the strongest candidates for conservation attention. Links that changed substantially
between assumptions are useful places for field surveys because better local data
could change the result. Agreement between scenarios is not complete proof because
the scenarios share many of the same input layers.}

% Survey implications follow directly from the uncertainty discussion.
\TODO{PLAIN-LANGUAGE GUIDE: there is no completed score that ranks individual survey
cells. Instead, explain how an organization could use the 13 uncertain links to choose
general survey areas, then adjust the order based on access, cost, safety, and local
knowledge. Link 17--24 is especially useful for checking because it is the only link
holding two parts of the three-neighbor network together, yet its mapped route was
highly unstable and had no overlap in 2020.}
